% Partie PFU (source : ../pfu/note-pfu.tex, corrigée selon ../revue/verif_is_pfu.md)
Le prélèvement forfaitaire unique (PFU) impose depuis 2018 les dividendes, les intérêts et les plus-values
mobilières à un taux unique, 31,4~\% en 2026, au lieu du barème progressif. La référence le garde tel qu'il
est :
\begin{itemize}\setlength{\itemsep}{0pt}
\item Quand le taux a monté en 2013, les actionnaires qui contrôlent leur société ont laissé les bénéfices
dedans. Quand il a baissé en 2018, ils les ont distribués.
\item Cette réaction porte sur le moment où l'on se verse le revenu plus que sur l'activité. Elle dépend sans
doute du taux qu'on attend pour demain : le principal service à rendre est de ne plus le changer.
\item La faiblesse du système n'est pas dans le taux du PFU, mais dans le circuit des holdings.
\end{itemize}

\subsection{Le PFU en 2026}

Le PFU additionne 12,8~\% d'impôt sur le revenu et 18,6~\% de prélèvements sociaux.
\begin{itemize}\setlength{\itemsep}{0pt}
\item Depuis la LFSS pour 2026, la CSG sur les revenus du capital est de 10,6~\%. Elle reste à 9,2~\% pour une
liste fermée de revenus (assurance-vie, épargne logement, revenus fonciers, plus-values immobilières) : deux
taux coexistent, sans justification économique.
\item Le contribuable peut opter pour le barème, globalement et, depuis 2026, de façon révocable. Sur leurs
revenus de 2024, 2,5~millions de foyers l'ont fait.
\item En haut de la distribution, la contribution exceptionnelle sur les hauts revenus (CEHR) ajoute 3 puis
4~points au-delà de 250\,000~€ de revenu (500\,000~€ pour un couple). La contribution différentielle (CDHR)
garantit un impôt d'au moins 20~\% du revenu, ce qui porte à 38,6~\% le taux d'un foyer qui vit de
dividendes ; elle est prorogée jusqu'à ce que le déficit du budget de l'État repasse sous 3~\% du PIB.
\end{itemize}

Le PFU ne s'apprécie qu'avec l'IS déjà payé par la société. Sur 100~€ de bénéfice, une société paie 25~€
d'IS (25,8~€ si son impôt dépasse 763\,000~€, avec la contribution sociale).
\begin{itemize}\setlength{\itemsep}{0pt}
\item Si elle distribue le reste, l'actionnaire paie 31,4~\%. Le bénéfice distribué supporte donc 48,6 à
49,1~\% au total, et jusqu'à 54,5~\% pour un foyer à la CDHR (figure~\ref{fig:taux}). C'est proche du taux
marginal supérieur sur un salaire (55,4~\% hors cotisations sociales).
\item Ce niveau est élevé en Europe. La France est sixième sur 35 pour le taux personnel supérieur sur les
dividendes (Tax Foundation, 2026, avec le taux de 34~\% de 2025) ; à cinq fois le salaire moyen, l'OCDE
trouvait en 2023 un taux moyen intégré de 49~\% que seuls le Danemark et l'Espagne dépassent.
\item Si la société verse son dividende à une holding de l'actionnaire qui le garde, la holding ne paie que
1,25~\% (régime mère-fille) et le bénéfice n'a supporté que 26,8~\%. C'est un report, pas un taux final :
distribué un jour, il paiera le PFU.
\end{itemize}
Imposer le bénéfice quand il sort des sociétés est d'ailleurs, en principe, ce qui gêne le moins la
croissance, puisque le capital circule d'une entreprise à l'autre sans impôt. Le report ne devient une
exonération que si la sortie n'a jamais lieu, ce que permet l'effacement au décès.

\figbloc{pfu_01_taux_cumule}{Impôt total sur 100~€ de bénéfice, selon ce qu'en fait l'actionnaire (2026)}{fig:taux}

\subsection{Ce que 2013 et 2018 ont appris}

Les dividendes ont suivi les deux réformes, en sens inverse :
\begin{itemize}\setlength{\itemsep}{0pt}
\item En 2013, l'option pour un prélèvement libératoire disparaît et tous les dividendes passent au barème.
Les foyers qui avaient opté pour le prélèvement libératoire déclarent 40~\% de dividendes en moins, sans
hausse de leurs salaires ; les entreprises contrôlées par des particuliers distribuent moins, accumulent des
actifs financiers et n'investissent pas moins (Bach et al., 2019).
\item Le taux marginal supérieur intégré n'était pourtant passé que de 57,6 à 60,1~\% : les actionnaires ont
pu juger la hausse provisoire et attendre, ou garder le bénéfice dans la société pour le sortir plus tard en
plus-value ou le transmettre (Bach et al., 2024).
\item En 2018, les dividendes remontent (+15~\% dès la première année). Le coût de la réforme tombe de 0,9 à
0,4~Md€ selon l'IPP (de 1,05 à 0,74~Md€, voire un gain dans l'hypothèse haute, selon l'Insee).
\item Mais 2013 et 2018 forment un seul épisode : les bénéfices retenus de 2013 à 2017 ont été distribués à
partir de 2018. L'« autofinancement » que France Stratégie (2023) déduit d'une part causale supposée, faute
de pouvoir la mesurer, ne vaut que pour les premières années, rattrapage compris.
\end{itemize}

Pour une société qui investit ses bénéfices, l'impôt sur le dividende ne change pas le choix d'investir :
100~€ distribués aujourd'hui rapportent 68,6~€ nets ; réinvestis à 5~\%, ils donneront 105~€ l'an prochain,
soit 72,0~€ nets. Le PFU prend la même fraction dans les deux cas : il change seulement la date et la forme
du revenu. Le raisonnement ne vaut pas pour une jeune entreprise qui lève des fonds propres à l'extérieur.
Les effets réels mesurés sont minces :
\begin{itemize}\setlength{\itemsep}{0pt}
\item après 2018, l'effet le plus net est une création d'entreprises un peu plus rapide dans les secteurs
exposés, qu'on ne sait pas séparer de celui de la baisse de l'IS (IPP, 2023) ;
\item la baisse américaine de 2003 n'a pas davantage modifié l'investissement (Yagan, 2015).
\end{itemize}

Sur le niveau du taux :
\begin{itemize}\setlength{\itemsep}{0pt}
\item Le taux qui maximise les recettes sur les revenus du capital, prélèvements sociaux compris, est estimé
à 43~\% en comptant la réaction des revenus du travail, à 57~\% sans elle (Lefebvre, Lehmann et Sicsic,
2025). Le PFU est en dessous, mais la marge est mince : avec l'élasticité estimée (0,77), un passage de 31,4
à 33~\% ne garde qu'un peu moins des deux tiers de son rendement mécanique.
\item Cette élasticité ne mesure pas une perte sèche : la réaction des dividendes consiste surtout à laisser le
bénéfice dans la société, ce qui reporte l'impôt sans détruire d'activité.
\item Les dividendes sont concentrés : 1~\% des foyers en déclarent 96~\% (2021). L'exil est faible (un point
d'imposition effective en plus fait partir à long terme 0,02 à 0,23~\% du centile aux plus hauts revenus du
capital ; Bach et al., 2025) et non mesuré au-delà.
\item La commission d'enquête de l'Assemblée sur les plus hauts patrimoines recommande de « maintenir la
stabilité » des taux (rapport n°~3056, 2026).
\end{itemize}

\subsection{Les arbitrages}

\textbf{Le taux.}
\begin{itemize}\setlength{\itemsep}{0pt}
\item Un PFU à 25~\% coûterait au plus 5,1~Md€ à comportements inchangés (un majorant, à cause des foyers qui
optent pour le barème) et 3,2 à 3,9~Md€ après réaction. Le gain irait d'abord au centile le plus aisé
(44~\% du gain du PFU selon l'estimation du Gouvernement en 2017).
\item Cette baisse ne rapporterait rien en activité : après 2018, aucun effet n'a été détecté sur
l'investissement ni sur les salaires des entreprises les plus exposées (France Stratégie, 2023). Elle
coûterait le prix des 3,7~Md€ à lever ailleurs, environ +0,3~Md€ de perte sèche (0,15 à 0,6).
\item La réaction des dividendes est une affaire de recette, pas d'efficacité : si une partie était une
évasion définitive que la baisse ferait cesser, la recette perdue serait moindre, sans gain d'activité.
\item Une hausse rapporterait peu : 1,2~Md€ nets pour 3~points, de 30 à 33~\%, selon le CAE (2025).
\end{itemize}

\textbf{Le report de l'apport-cession.} Un entrepreneur qui apporte ses titres à une holding avant de les
vendre place sa plus-value en report : elle n'est imposée que s'il vend les titres de la holding, ou si la
holding vend dans les trois ans sans réinvestir au moins 70~\% du produit. Au décès, le report devient une
exonération définitive ; en cas de donation, il s'efface si le donataire ne vend pas pendant six ans (onze
dans certains cas). D'après le rapport n°~3056 :
\begin{itemize}\setlength{\itemsep}{0pt}
\item 33,6~Md€ de plus-values ont été placées en report en 2023 et 22,2~Md€ en 2024, plus que les plus-values
imposées au PFU en 2024 (18,8~Md€ d'après les déclarations) ;
\item plus des deux tiers des entrées de 2023 (23,5~Md€) se concentrent sur 84 foyers ;
\item le stock atteint 153~Md€ en 2023 et 173~Md€ en 2024, et seuls 0,4 à 1,5~Md€ de reports ont expiré
chaque année de 2022 à 2024.
\end{itemize}
Trois traitements sont possibles :
\begin{itemize}\setlength{\itemsep}{0pt}
\item \emph{L'expiration à toute transmission} (PLF 2027, art.~5) : 0,5~Md€ annoncés pour 2027 et, selon
notre calcul, au moins 2~Md€ par an en régime de croisière. Mais elle taxe une structure plutôt qu'un gain
réalisé : une holding familiale transmise sous pacte Dutreil paierait sans avoir rien vendu, et l'étalement
sur cinq ans ne vise que l'impôt sur le revenu. Elle alourdit aussi de 90 à 100~M€ chacun la facture des plus gros apporteurs.
\item \emph{L'expiration seulement quand la holding vend.}
\item \emph{La transmission du report aux héritiers sans effacement} : la plus-value est imposée quand le
capital sort réellement des sociétés, sans impôt sans liquidité ni hausse de facture à la transmission. Elle
rapporte moins et plus tard, et sa solidité juridique reste à expertiser.
\end{itemize}

\textbf{Les autres choix.}
\begin{itemize}\setlength{\itemsep}{0pt}
\item Dans l'impôt fusionné, la CEHR disparaît pour les revenus du barème. La garder sur les revenus du
capital, sous la forme d'une majoration de 3 puis 4~points du PFU, maintient le taux actuel (tant que la CDHR
existe, un foyer qui vit de dividendes, déjà au plancher de 20~\%, n'y gagnerait de toute façon rien). Durcir
la CDHR, comme le propose la commission d'enquête, alourdirait la facture des mêmes foyers.
\item Aligner les deux taux de prélèvements sociaux à 18,6~\% rapporterait de l'ordre de 1,3~Md€, payés
surtout par la classe moyenne épargnante.
\item Un prélèvement modéré sur la trésorerie passive des holdings (l'Irlande taxe à 20~\% les revenus passifs
que les sociétés fermées ne distribuent pas dans les dix-huit mois) viserait le bon problème, mais
alourdirait lui aussi la facture des grandes fortunes. Un impôt de 25~\% sur les revenus non réinvestis des
holdings risque de heurter la directive mère-fille, selon l'administration fiscale, et, s'il est établi chez
l'associé, la jurisprudence constitutionnelle (imposer des bénéfices qu'il n'a pas perçus) ; il taxerait le
capital avant qu'il sorte des sociétés.
\item Pour une détention longue, la correction de principe est d'imposer le gain réel, en indexant le prix
d'acquisition sur l'inflation, plutôt que de rétablir des abattements pour durée de détention. Ces
abattements avaient compliqué l'impôt de 2013 à 2017.
\end{itemize}
